% sec7a.tex -- writer W4. Section 7 heading, 7.1 (constants), 7.2 (run 1, tangency, forward arc, virtual labels).
\section{The front at touch-down}\label{sec:touch}

In this section $P$ is design T (Definition~\ref{def:T}) and nothing after $\ttd$ is considered. The fixed data are
\[
\vbar=5.8\cdot10^{-3},\quad\eta=0.11639,\quad\epsbar=10^{-3},\quad\xi_1=1.16171,\quad\lambda_0=0.0337254,\quad G_{\mathrm T}=0.030573,
\]
and throughout $0<v\le\vbar$, $0<\eps\le\epsbar$. We put $\mub:=\vbar/(1-\vbar)$, $\muT:=3v/(1+v)$, $\muTb:=3\vbar/(1+\vbar)$; so
$\mu\le\mub$ and $\muT\le\muTb$. The notation of Section~\ref{sec:designT} is used throughout ($t_1$, $\alpha_0$, $\rho_0$, $\beta^*$,
$\nhat$, $\ta$, $\ttd$, $\tc$, $X_0$, $\omega$; run-1 data $\rho_1,\beta_1,\tau_1,\tau_\alpha,y(\alpha,t_1),h_1$; run-2 data
$\sr,\sbb,\tau_{\sbb},C,\sa,\sfc,\sA$; $I_\eta$, $\xi$).

\subsection{The constants of the touch-down}\label{sec:tdconst}

In the next definition every formula is evaluated with $v:=\vbar$, $\mu:=\mub$, $\muT:=\muTb$, $\eps:=\epsbar$; Lemma~\ref{lem:mono}
explains why this is legitimate. The numerical values are listed in Table~\ref{tab:td}.

\begin{definition}[touch-down constants]\label{def:tdconst}
{\small
\begin{itemize}\setlength\itemsep{2pt}
\item \emph{Basic.} $\bar t_1:=2+\mub$; $\gb:=(1-\vbar^2)^{1/2}((1+\vbar)/(1-\vbar))^{\vbar/2}$ $(=\gmax$ at $v=\vbar)$;
$\underline\alpha:=\arccos(\frac12+\vbar)$; $\Delta_{\mathrm a}:=\pi/3-\underline\alpha$; $d_1:=\vbar\bar t_1+\mub$;
$e_\rho:=\vbar\bar t_1+\Delta_{\mathrm a}+\mub$; $T_2:=\bar t_1+(\sqrt3+e_\rho)/(1-\vbar)$; $d_{\mathrm s}:=\vbar T_2$;
$d_P:=\max(\mub,d_{\mathrm s}+\Delta_{\mathrm a})$.
\item \emph{Run 1.} $\rho_{\mathrm{mid}}:=\sin\frac12$; for $\rho>0$: $W_1(\rho):=e^{\vbar\bar t_1/\rho}$, $\omega_1(\rho):=\vbar\bar t_1/(2\rho)$,
$\vk(\rho):=(2/\rho+1/\rho^2)\vbar\bar t_1e^{2\vbar\bar t_1/\rho}$; $k_{\mathrm{arm}}:=(1+\vbar)/(\cos\frac12-\vbar)+2\vbar/(1-\vbar^2)$;
$d_{\mathrm{arm}}:=(1+\vbar)/((1-\vbar)(\cos\frac12-\vbar))+\vbar/(1-\vbar^2)$;
$e_1:=(2-\cos\frac12)(1+\tan^2\frac14)+(1+\vbar)\tan(\frac14)^{1-\vbar}$, $\mathcal E_1:=\max(1,e_1/(1-\vbar))$,
$W_{\mathrm{arm}}:=\mathcal E_1^{\vbar}/(1-\vbar)$; $k_{\mathrm{all}}:=\max(k_{\mathrm{arm}},1+\vk(\rho_{\mathrm{mid}}))$,
$W_{\mathrm{all}}:=\max(W_{\mathrm{arm}},W_1(\rho_{\mathrm{mid}})(\pi/2)^{2\vbar})$.
\item \emph{Window.} $\theta_L:=\underline\alpha-\eta$, $\theta_R:=\pi/3+\eta$; $R_0(\theta):=(5-4\cos\theta)^{1/2}=|e^{i\theta}-2|$,
$R_L:=R_0(\theta_L)$, $R_R:=R_0(\theta_R)$; $\sr_{\min}:=R_L-d_1$, $\sr_{\max}:=R_R+d_1$; $\rho_{\mathrm w}:=\sin\theta_L$,
$W_{\mathrm w}:=W_1(\rho_{\mathrm w})$, $\omega_{\mathrm w}:=\omega_1(\rho_{\mathrm w})$, $\vk_{\mathrm w}:=\vk(\rho_{\mathrm w})$;
$\iota_{\mathrm w}:=1+\vbar\bar t_1/\sin\theta_L$; $q_{\min}:=\sin(\underline\alpha+\eta/2)/R_R$, $q_{\max}:=\sin(\pi/3-\eta/2)/R_L$;
$\ell_{\min}:=4q_{\min}/(1+d_1/R_L)$, $\ell_{\max}:=2\bar t_1\iota_{\mathrm w}q_{\max}/(1-d_1/R_L)$;
$\Lambda_{\min}:=4\sin(\underline\alpha-\eta/2)/\sr_{\max}^2$, $\Lambda_{\max}:=2\bar t_1\iota_{\mathrm w}\sin(\pi/3+\eta/2)/\sr_{\min}^2$;
$m_0:=\max(|1-2\cos\theta_L|,|1-2\cos\theta_R|)/R_L$, $\omega_{\mathrm p}:=\arcsin(d_1/R_L)$,
$\angle_{\mathrm w}:=\arcsin m_0+\omega_{\mathrm w}+\omega_{\mathrm p}$, $\sA_{\min}:=\cos\angle_{\mathrm w}$, $\sA_{\max}:=W_{\mathrm w}$;
$b_{\mathrm{lo}}:=\Lambda_{\min}(1-\eta^2/48)/4$, $b_{\mathrm{hi}}:=\Lambda_{\max}/4$, $\bar b:=b_{\mathrm{hi}}\eta^2$,
$e_\tau:=\tan(\bar b/2)/(\bar b/2)-1$, $\beta_2:=b_{\mathrm{hi}}(1+e_\tau)/2$;
$C_{\mathrm{hi}}:=2\sr_{\max}\beta_2^{1-\vbar}$, $C_{\mathrm{lo}}:=2\sr_{\min}(b_{\mathrm{lo}}/2)/(1+\tan^2(\bar b/2))$;
$\mathrm{aux}:=\ell_{\max}\bar b\eta/2$, $\sa_{\min}:=\sA_{\min}\cos\bar b-\mathrm{aux}$, $\sa_{\max}:=\sA_{\max}(1+\vbar)+\mathrm{aux}$;
$L_C:=2\Lambda_{\max}/\Lambda_{\min}+\eta\sA_{\max}/\sr_{\min}$,
$L_{\mathrm{lo}}:=2(\cos\bar b-\vbar)(\Lambda_{\min}/\Lambda_{\max})(1-\eta^2/32)$;
$S_1:=\ell_{\max}(\sr_{\max}\sa_{\max}L_C+(1+\vbar)\eta)/(2(1-\vbar^2)\sr_{\min})$.
\item \emph{Forward arc.} $h_{\mathrm f}:=\sA_{\min}-\sr_{\max}b_{\mathrm{hi}}^2\eta^3/2$, $\rho_{\mathrm c}:=h_{\mathrm f}/(1+\vbar)$,
$Q:=(1+\vbar)\sr_{\max}/h_{\mathrm f}$, $y_{\mathrm f}:=2\sr_{\max}\beta_2$;
$Y_1:=(\ell_{\max}/2+\sA_{\max}b_{\mathrm{hi}}(1+e_\tau)\eta)/\sa_{\min}$, $\bar Y:=Y_1\eta^{1-\mub}Q^{\mub}$,
$\theta_{\mathrm f}:=\max((1+\vbar)Y_1,\vbar y_{\mathrm f}/\rho_{\mathrm c})/(1-\vbar-\vbar\bar Y)$,
$W_{\mathrm f}:=\sa_{\max}[(1+\vbar+\vbar\bar Y)^2+((1+\vbar)\bar Y+\vbar)^2]^{1/2}/(1-\vbar^2)$;
$\delta_{\mathrm{fw}}:=\vbar/(3(1-\vbar)^2\xi_1)$, $\tau_{\mathrm f}:=\beta_2Q^{1+\mub}$, $\bar\tau_{\mathrm f}:=\tau_{\mathrm f}\eta^{1-\mub}$,
$\rho_{\mathrm u}:=((1+\vbar)\sA_{\max}+\delta_{\mathrm{fw}})/(1-2\bar\tau_{\mathrm f}^2)$,
$S_2:=(2\tau_{\mathrm f}+\theta_{\mathrm f}Q^{\mub})(\rho_{\mathrm u}\sa_{\max}L_C+1+\vbar)/((1-\vbar^2)\rho_{\mathrm c})$,
$k_{\mathrm f}:=Q^{\mub}(1+\vk_{\mathrm w}+S_1\vbar)+S_2\vbar$.
\item \emph{Virtual labels.} $\theta_{\mathrm v}:=Q^{\mub}e^{\mub}(\theta_{\mathrm f}+\mub e^{2\mub}y_{\mathrm f}\xi_1)$,
$W_{\mathrm v}:=W_{\mathrm f}e^{2\mub}$, $u_{\mathrm v}:=W_{\mathrm v}\xi_1$.
\item \emph{Tail.} $\xi_-:=C_{\mathrm{hi}}^{-1/2}$, $\xi_+:=(\gb/C_{\mathrm{lo}})^{1/(2-2\vbar)}$, $\bar\xi_{\mathrm e}:=\xi_+\epsbar^{1/2}$;
$\delta_{\mathrm{tl}}:=(1+\vbar)\vbar/(3(1-\vbar)\xi_-)$, $h_{\mathrm t}:=\sA_{\min}-\delta_{\mathrm{tl}}$,
$h_{\mathrm T}:=(1+\vbar)\sA_{\max}+\sr_{\max}b_{\mathrm{hi}}^2\eta^3/2$; $m_\tau:=(2h_{\mathrm t}/C_{\mathrm{hi}})\eta^{-(1-2\vbar)}$,
$\tau_{\mathrm t}:=m_\tau^{1/(1+\vbar)}$, $e_{\mathrm t}:=(1+\vbar)/(m_\tau\tau_{\mathrm t}^{1-\vbar})$;
$y_{\mathrm c}:=2(1+e_{\mathrm t})h_{\mathrm T}/(1-\vbar)$, $y_{\mathrm T}:=C_{\mathrm{hi}}^{1/(1+\vbar)}\max(1,y_{\mathrm c})^{\vbar}$,
$T_{\mathrm t}:=(\max(1,y_{\mathrm c}/C_{\mathrm{lo}})\cdot2/b_{\mathrm{lo}})^{\vbar}$;
$Y_2:=T_{\mathrm t}\ell_{\max}/(2\sa_{\min})$, $\bar Y_2:=Y_2\eta^{1-\muTb}$,
$\theta_{\mathrm t}:=Y_2+\vbar y_{\mathrm T}/h_{\mathrm t}$, $z_{\mathrm T}:=(h_{\mathrm T}+\vbar y_{\mathrm T}\eta^{1-\muTb})/(1-\vbar)$,
$\rho_{\mathrm T}:=z_{\mathrm T}(1+(y_{\mathrm T}\eta^{1-\muTb}/h_{\mathrm t})^2)^{1/2}$;
$W_{\mathrm t}:=\sa_{\max}[(1+\vbar+\vbar\bar Y_2)^2+((1+\vbar)\bar Y_2+\vbar)^2]^{1/2}/(1-\vbar^2)$;
$S_3:=(y_{\mathrm T}/h_{\mathrm t}+\theta_{\mathrm t})(\rho_{\mathrm T}\sa_{\max}L_C+1+\vbar)/((1-\vbar^2)h_{\mathrm t})$,
$k_{\mathrm t}:=T_{\mathrm t}(1+\vk_{\mathrm w}+S_1\vbar)+S_3\vbar$; $\xi_2:=1/(3(1-\vbar)\sA_{\min})$;
$h_{\mathrm a}:=1-\sA_{\max}/2-\sr_{\max}b_{\mathrm{hi}}^2\epsbar^3/32$, $Q_{\mathrm a}:=(1+\vbar)\sr_{\max}/h_{\mathrm a}$,
$\tau_{\mathrm a}:=\beta_2Q_{\mathrm a}^{1+\mub}\epsbar^{1-\mub}/4$.
\item \emph{Bulk and hole.} $e_\beta:=\arcsin(e_\rho/\sqrt3)$, $b_{\mathrm l}:=b_{\mathrm{lo}}\eta^2$, $y_{\mathrm b}:=\sin b_{\mathrm l}$,
$H_{\mathrm b}:=3+d_1+(1+\vbar)(\sqrt3+e_\rho)/(1-\vbar)$, $T_{\mathrm b}:=2H_{\mathrm b}/(y_{\mathrm b}\tan(b_{\mathrm l}/2))$,
$S_{\mathrm b}:=(3+\vbar)/(2(1-\vbar)y_{\mathrm b})$, $k_{\mathrm b}:=T_{\mathrm b}^{\vbar}(k_{\mathrm{all}}+S_{\mathrm b}\vbar)+S_{\mathrm b}\vbar$,
$W_{\mathrm b}:=((1+\vbar)/(1-\vbar))T_{\mathrm b}^{\vbar}$; $r_{\mathrm B}:=0.005+\vbar\bar t_1$,
$\rho_{\mathrm h}:=\frac12-d_P-r_{\mathrm B}-\vbar T_2$.
\item \emph{Assembly.} $W_x:=\max(W_{\mathrm v},W_{\mathrm t},W_{\mathrm b}W_{\mathrm{all}})$,
$U_{\max}:=(W_{\mathrm v}+W_{\mathrm t})\eta+W_{\mathrm b}W_{\mathrm{all}}(2\pi-2\eta)/(1-2\vbar)$;
$\hat L:=(8/\lambda_0)^{2\vbar}/(1-2\vbar)$, $\Delta_0:=\lambda_0/(W_x\hat L)$, $c_{\mathrm{sep}}:=2\sin(\Delta_0/2)$,
$\kc:=\min(\cos(\theta_{\mathrm v}\Delta_0^{1-\mub}),h_{\mathrm t}/W_{\mathrm t})$;
$k_*:=\max(k_{\mathrm f}W_{\mathrm v}^{\mub},k_{\mathrm t}W_{\mathrm t}^{\muTb},k_{\mathrm b})$, $c_{\mathrm{far}}:=k_*/\sqrt{12}$,
$c_{\mathrm{near}}:=\theta_{\mathrm v}((12u_{\mathrm v}^2+6u_{\mathrm v}+1)/(3-2\mub))^{1/2}$,
$c_{\mathrm{tail}}:=\theta_{\mathrm t}(19/(3-2\muTb))^{1/2}$, $c_N:=\max(c_{\mathrm{far}},c_{\mathrm{near}},c_{\mathrm{tail}})$;
$c_y:=y_{\mathrm T}/h_{\mathrm t}^2$; $c_{T4}:=\theta_{\mathrm t}/h_{\mathrm t}$.
\end{itemize}}
\end{definition}

\noindent\textbf{Side conditions.}\namedlabel{hyp:SC}{(SC)} {\small
\hyptag{SC1} $\sqrt3-e_\rho>\epsbar$. \hyptag{SC2} $\cos\frac12>\vbar$. \hyptag{SC3} $\theta_L\ge\frac12$. \hyptag{SC4} $\theta_R\le\pi/2$.
\hyptag{SC5} $\eta\le\pi/3$. \hyptag{SC6} $\cos\theta_L<4/5$. \hyptag{SC7} $\cot(\underline\alpha-\eta/2)/2<s(\theta_R)$, where
$s(\theta):=2\sin\theta/(5-4\cos\theta)$. \hyptag{SC8} $d_1<R_L$. \hyptag{SC9} $\angle_{\mathrm w}<\pi/2$. \hyptag{SC10} $\bar b<\pi/2$.
\hyptag{SC11} $\beta_2<1$. \hyptag{SC12} $b_{\mathrm{lo}}<2$. \hyptag{SC13} $\sa_{\min}>0$. \hyptag{SC14} $\cos\bar b>\vbar$.
\hyptag{SC15} $L_{\mathrm{lo}}>\eta\sA_{\max}/\sr_{\min}$. \hyptag{SC16} $0<h_{\mathrm f}\le1-\vbar^2$. \hyptag{SC17} $\rho_{\mathrm c}\xi_1>1$.
\hyptag{SC18} $Q\ge1$. \hyptag{SC19} $1-\vbar-\vbar\bar Y>0$. \hyptag{SC20} $1-2\bar\tau_{\mathrm f}^2>\vbar$. \hyptag{SC21} $C_{\mathrm{hi}}<1$.
\hyptag{SC22} $C_{\mathrm{lo}}\eta^2>\epsbar\gb$. \hyptag{SC23} $\gb>C_{\mathrm{lo}}$. \hyptag{SC24} $\bar\xi_{\mathrm e}<\eta$.
\hyptag{SC25} $h_{\mathrm t}>0$. \hyptag{SC26} $m_\tau>1$. \hyptag{SC27} $\xi_2\vbar\le\frac12$. \hyptag{SC28} $\xi_-\ge\frac12$.
\hyptag{SC29} $h_{\mathrm a}>0$. \hyptag{SC30} $1-2\tau_{\mathrm a}^2>0$. \hyptag{SC31} $\sA_{\min}>4\tau_{\mathrm a}^2$.
\hyptag{SC32} $b_{\mathrm l}\le\pi/6-e_\beta$. \hyptag{SC33} $\pi/3+2e_\beta\le\pi-b_{\mathrm l}$. \hyptag{SC34} $T_{\mathrm b}\ge1$.
\hyptag{SC35} $\vbar\le\frac1{10}$. \hyptag{SC36} $\rho_{\mathrm h}>0$. \hyptag{SC37} $2\gb\epsbar<\theta_L$ and $\gb\epsbar\le1$.
\hyptag{SC38} $2-2d_{\mathrm s}\ge1$. \hyptag{SC39} $2\lambda_0+\min(\bar\xi_{\mathrm e},\lambda_0)\le\eta$. \hyptag{SC40}
$2u_{\mathrm v}\epsbar\le\eta$. \hyptag{SC41} $W_x\hat L\le2$. \hyptag{SC42} $\Delta_0\le\eta$.
\hyptag{SC43} $c_{\mathrm{sep}}>2\epsbar$. \hyptag{SC44} $\theta_{\mathrm v}\Delta_0^{1-\mub}\le\pi/2$.
\hyptag{SC45} $k_*\lambda_0^{1-\muTb}\le\pi$. \hyptag{SC46} $u_{\mathrm v}\ge\frac12$.
\hyptag{SC47} $G_{\mathrm T}>0$ and $\bar\xi_{\mathrm e}+G_{\mathrm T}\le\eta$. \hyptag{SC48} $z_{\mathrm T}\bar\xi_{\mathrm e}+\epsbar<1$ and
$G_{\mathrm T}+z_{\mathrm T}\bar\xi_{\mathrm e}\le1$.}

\medskip
Each side condition is one inequality between explicitly computable numbers; all hold (checked in Appendix~\ref{app:const}).

\begin{lemma}[monotonicity]\label{lem:mono}
Assume \hyptag{SC1}--\hyptag{SC48}, and let $0<v\le\vbar$ and $0<\eps\le\epsbar$. Every constant of
Definition~\ref{def:tdconst} other than $\kc$ is the value at $(\vbar,\epsbar)$ of an
expression $E(v,\eps)$ (replace $\vbar,\mub,\muTb,\epsbar$ by $v,\mu,\muT,\eps$) which is non-decreasing in $v\in(0,\vbar]$ and in
$\eps\in(0,\epsbar]$ (an \emph{upper} constant) or non-increasing in both (a \emph{lower} constant). The lower constants are
$\underline\alpha$, $\rho_{\mathrm w}$, $\sr_{\min}$, $\ell_{\min}$, $\Lambda_{\min}$, $\sA_{\min}$, $b_{\mathrm{lo}}$, $C_{\mathrm{lo}}$,
$\sa_{\min}$, $L_{\mathrm{lo}}$, $h_{\mathrm f}$, $\rho_{\mathrm c}$, $\xi_-$, $h_{\mathrm t}$, $m_\tau$, $\tau_{\mathrm t}$, $h_{\mathrm a}$,
$b_{\mathrm l}$, $y_{\mathrm b}$, $\rho_{\mathrm h}$, $\Delta_0$, $c_{\mathrm{sep}}$, $\theta_L$, $R_L$, $q_{\min}$; all others except $\kc$
are upper constants.
Consequently, wherever a proof uses a constant as an upper (lower) bound for a quantity formed with the actual $v,\mu,\muT,\eps$, the
constant is $\ge$ ($\le$) the value of that expression; and every side condition other than \hyptag{SC44} implies the same
inequality at $(v,\eps)$.
\end{lemma}
\begin{proof}
On $(0,\vbar]\subset(0,\frac1{10}]$ the following are increasing in $v$: $v$, $\mu$, $\muT$, $1+v$, $1/(1-v)$, $1/(1-v^2)$,
$t_1=2+\mu$, $vt_1$ and $\gmax(v)$ ($d\log\gmax/dv=\operatorname{artanh}v>0$, from the closed form in Proposition~\ref{prop:runs});
$1-v$, $\cos\frac12-v$, $\cos\bar b-v$ and $\arccos(\frac12+v)$ are decreasing. Powers are handled by three rules:
(1) $X^{p(v)}$ with $X\ge1$ and $p$ non-decreasing is non-decreasing: $Q^\mu$, $Q^{1+\mu}$, $Q_{\mathrm a}^{1+\mu}$, $\mathcal E_1^v$,
$(\pi/2)^{2v}$, $T_{\mathrm b}^v$, $(8/\lambda_0)^{2v}$, $W_{\mathrm v}^\mu$, $W_{\mathrm t}^{\muT}$,
$(\max(1,y_{\mathrm c}/C_{\mathrm{lo}})2/b_{\mathrm{lo}})^v$ (base $\ge1$ by \hyptag{SC12}), $\max(1,y_{\mathrm c})^v$,
$(\gb/C_{\mathrm{lo}})^{1/(2-2v)}$ (base $>1$ by \hyptag{SC23}); here $Q\ge1$ is \hyptag{SC18}, $Q_{\mathrm a}\ge(1+v)\sr_{\max}\ge1$ as
$h_{\mathrm a}\le1$ and $R_0\ge1$, and $W_{\mathrm v},W_{\mathrm t}\ge\sa_{\max}\ge\sA_{\max}=W_{\mathrm w}\ge1$;
(2) $X^{p(v)}$ with $0<X\le1$ and $p$ non-increasing is non-decreasing: $\beta_2^{1-v}$ (\hyptag{SC11}), $C_{\mathrm{hi}}^{1/(1+v)}$
(\hyptag{SC21}), $\eta^{1-\mu}$, $\eta^{1-\muT}$, $\epsbar^{1-\mu}$, $\tan(\frac14)^{1-v}$, $\lambda_0^{1-\muT}$;
(3) $X^{-p(v)}$ with $0<X<1$ and $p$ non-increasing is non-increasing: $\eta^{-(1-2v)}$ in $m_\tau$; then $m_\tau^{1/(1+v)}$ with
$m_\tau>1$ is a lower constant by \hyptag{SC26}, and $\tau^{v-1}\le\tau_{\mathrm t}^{v-1}\le\tau_{\mathrm t}^{\vbar-1}$ for
$\tau\ge\tau_{\mathrm t}\ge1$ (in $e_{\mathrm t}$). The constant $\kc$ and \hyptag{SC44} enter the proofs only as fixed numbers
(in (T2), $\xi<\Delta_0$ refers to the fixed $\Delta_0$), so no monotonicity is needed for $\kc$ or $\Delta_0^{1-\mu}$. Every constant is built from these blocks, from the fixed data and from earlier
constants by sums, products, quotients of positive quantities, $\max$, $\min$, and the monotone functions $\exp$, $\sqrt\cdot$,
$\arcsin$ on $[0,1]$, $\sin$ and $\tan$ on $[0,\pi/2)$, $\cos$ on $[0,\pi/2]$ and $\cot$ on $(0,\pi/2)$; the arguments stay in these
ranges by \hyptag{SC1}, \hyptag{SC8}--\hyptag{SC10} and $m_0\le1$ (checked in Appendix~\ref{app:const}). Going through Definition~\ref{def:tdconst} in order, each operation combines upper
into upper and lower into lower constants in the direction stated (upper minus lower is upper, upper over lower is upper, etc.).
Three further steps use side conditions: $C_{\mathrm{lo}}$ uses $(b_{\mathrm{lo}}\xi^2/2)^{1-v}\ge(b_{\mathrm{lo}}/2)\xi^{2-2v}$
(\hyptag{SC12}); $\xi_-=C_{\mathrm{hi}}^{-1/2}\le C_{\mathrm{hi}}^{-1/(2-2v)}$ (\hyptag{SC21}); $\bar Y$ and $\theta_{\mathrm f}Q^{\mub}$
use $Q\ge1$. In $\eps$: the constants containing $\epsbar$ are $\bar\xi_{\mathrm e}$, $h_{\mathrm a}$, $Q_{\mathrm a}$, $\tau_{\mathrm a}$,
all monotone in the stated direction. Each side condition other than \hyptag{SC12}, the second part of \hyptag{SC16}, \hyptag{SC18}, \hyptag{SC23}, \hyptag{SC32}--\hyptag{SC34}, \hyptag{SC42} and \hyptag{SC46} compares an upper with a lower expression (or a number) in the direction that is preserved when $(\vbar,\epsbar)$ is replaced by $(v,\eps)$. The exceptions hold for every $v\in(0,\vbar]$ and $\eps\in(0,\epsbar]$ directly. Since $R_0\ge1$ we have $\sr_{\max}\ge1$, hence $b_{\mathrm{lo}}\le\Lambda_{\min}/4\le1$ \hyptag{SC12}, $b_{\mathrm l}\le\eta^2$, and $C_{\mathrm{lo}}\le\sr_{\min}b_{\mathrm{lo}}\le1/\sr_{\max}\le1/R_R<1\le\gb$ \hyptag{SC23}. Since $h_{\mathrm f}\le\sA_{\min}\le1$, $Q\ge\sr_{\max}\ge1$ \hyptag{SC18}. The bound $\eta^2+2e_\beta<\pi/6$ gives \hyptag{SC32} and \hyptag{SC33}. Since $y_{\mathrm b}\tan(b_{\mathrm l}/2)<1<H_{\mathrm b}$, $T_{\mathrm b}>1$ \hyptag{SC34}. Since $W_x,\hat L\ge1$, $\Delta_0\le\lambda_0<\eta$ \hyptag{SC42}. Since $W_{\mathrm v}\ge1$, $u_{\mathrm v}\ge\xi_1>\frac12$ \hyptag{SC46}. Finally, the lower constant $h_{\mathrm f}$ is at most its value $0.99152$ at $v=0$, which is $<1-\vbar^2\le1-v^2$ (second part of \hyptag{SC16}).
\end{proof}

From now on Lemma~\ref{lem:mono} is used tacitly: ``by $v\le\vbar$'' means that a $v$-expression is bounded by its constant.

\subsection{Run 1, tangency and the forward arc}\label{sec:fwd}

\begin{lemma}[run 1]\label{lem:run1}
\begin{enumerate}
\item[(a)] Let $\as\in(0,\arccos v)$ be given by $\sg(\as)=\eps\gmax$. A true candidate $\alpha\in(-\pi,\pi]$ is seen during
$[0,t_1]$ iff $|\alpha|\le\as$. For $\alpha\in[\as,2\pi-\as]$, $\rho_\alpha\ge\eps$ on $[0,t_1]$, so virtual $=$ true on $[0,t_1]$.
Moreover $\as\le2\gb\eps<\theta_L$, and every candidate, true or virtual, satisfies $|S(t)-e^{i\alpha}|\le vt$.
\item[(b)] \emph{Middle.} For $\alpha\in[\frac12,2\pi-\frac12]$, $\rho_\alpha\ge\rho_{\mathrm{mid}}$ on $[0,t_1]$; for
$\alpha\in[\frac12,\frac\pi2]$ (and the mirror labels), $\rho_\alpha\ge\sin\alpha$ on $[0,t_1]$. If $\rho_\alpha\ge\rho$ on $[0,t_1]$,
then at $t_1$: $1\le w_1\le W_1(\rho)$, $|\arg\partial_\alpha S_\alpha(t_1)-(\alpha+\pi/2)|\le\omega_1(\rho)$, $|k_1-1|\le\vk(\rho)$.
\item[(c)] \emph{Arm.} For $\alpha\in(0,\frac12]$: $\rho_1\ge1+v$, $\tau_1>1$, $y(\alpha,t_1)\ge\sin\alpha$; $|k_1|\le k_{\mathrm{arm}}$;
the tangent at $t_1$ is $e^{i(\pi/2+\delta)}$ with $0\le\delta\le\tan\delta\le d_{\mathrm{arm}}\,y(\alpha,t_1)$;
$w_1\le W_{\mathrm{arm}}(2/\alpha)^{2v}$; and $y(\alpha,t_1)\le2\mathcal E_1\tau_\alpha^{(1-v)/(1+v)}$.
\item[(d)] For all $\alpha\in(0,2\pi)$: $|k_1|\le k_{\mathrm{all}}$, $1\le w_1\le W_{\mathrm{all}}(2/\circab{\alpha})^{2v}$ and
$\rho_1\ge1+v$, where $\circab\alpha:=\min(\alpha,2\pi-\alpha)$.
\item[(e)] For every $\alpha\in(0,2\pi)$ the true candidate $S_\alpha$ exists on $[0,\ttd]$ (it does not meet $P$). The candidate
$\alpha=0$ is caught at $1/(1-v)<t_1$.
\end{enumerate}
In (b)--(d) the true front is meant, also at the labels $0<\circab\alpha<\as$, which have been seen but keep moving by the true law.
\end{lemma}
\begin{proof}
(a) By Lemma~\ref{lem:coord}(b),(c), $\tau$ increases from $\tau_\alpha$ to $\tau_1$, and $h_1<0$ means $\cos\beta_1<-v$, so
$\tau_1^2>(1+v)/(1-v)>(1-v)/(1+v)$. If $0<\alpha\le\arccos v$ the minimum point of $\rho$ is passed, and
$\min_{[0,t_1]}\rho=\sg(\alpha)/\gmax$, which is $\le\eps$ iff $\alpha\le\as$ ($\sg$ is increasing on $(0,\arccos v]$); if
$\alpha\in(\arccos v,\pi)$, $\rho$ is increasing and $\min\rho=1$. The label $\pi$ has $\beta\equiv\pi$ and $\rho\ge1$; the lower half is
the mirror image; $\alpha=0$ is caught (see (e)). Virtual $=$ true by Lemma~\ref{lem:truefront}. With $\tau_{\mathrm s}:=\tan(\as/2)\le1$:
$\tau_{\mathrm s}^{1-v}=\eps\gmax(1+\tau_{\mathrm s}^2)/2\le\eps\gmax$, so $\as\le2\tau_{\mathrm s}\le2(\eps\gmax)^{1/(1-v)}\le2\eps\gmax\le2\gb\eps$,
because $\eps\gmax\le\gb\epsbar\le1$; and $2\gb\eps<\theta_L$ (\hyptag{SC37}). Every candidate has speed $\le v$ (true: by definition;
virtual: $|F|\le1$, Lemma~\ref{lem:field}(i)).

(b) For $\alpha\in[\frac12,\frac\pi2]$: $y=C\tau^v$ is non-decreasing, so $\rho\ge y\ge\sin\alpha\ge\sin\frac12$. For
$\alpha\in[\frac\pi2,\pi]$: $X:=\operatorname{Re}S_\alpha$ has $X'=v(X-t)/\rho$ and $(X-t)'\le v-1<0$, $X(0)=\cos\alpha\le0$; so $X-t\le X\le\cos\alpha\le0$,
and with $y\ge\sin\alpha$, $\rho^2\ge\cos^2\alpha+\sin^2\alpha=1$. Mirror for the lower half. By Lemma~\ref{lem:truefront}, $(\log w)_t\in[0,v/\rho]$,
so $1\le w_1\le e^{vt_1/\rho}$; $|\psi_t|\le v/(2\rho)$, and the tangent at $t=0$ is $e^{i(\alpha+\pi/2)}$; and, as
$|1-3\sin^2\!\phi|\le2$ and $|\sin\phi(3\cos^2\!\phi-1)|\le1$ (the maximum of $|s(2-3s^2)|$ on $[-1,1]$ is $1$),
$\frac d{dt}|k-1|\le(2v/\rho)|k-1|+2v/\rho+v/\rho^2$, whence by Gronwall ($k(0)=1$) $|k_1-1|\le(2/\rho+1/\rho^2)vt_1e^{2vt_1/\rho}$. Since
$vt_1\le\vbar\bar t_1$, these are the bounds claimed.

(c) $h_1=-(1-v^2)-(1-\cos\alpha)<0$, so $\cos\beta_1<-v$, $\beta_1>\pi/2$, $\tau_1>1$; and $|h_1|=\rho_1|\cos\beta_1+v|\le\rho_1(1-v)$ gives
$\rho_1\ge1+v$. $y=C\tau^v$ is non-decreasing and $y(0)=\sin\alpha$.

\emph{Curvature.} Apply Lemma~\ref{lem:transport}(c) at $t=0$, where $w=1$, $k=1$, $\rho=1$, $\psi=\alpha+\pi/2$, $\phi=-\pi/2$ and
$\sa\,\partial_\alpha\log C=-\sin\alpha\cdot(\cos\alpha-v)/\sin\alpha$: there
$\TS=-\frac v{1-v^2}[-(\cos\alpha-v)-v-\cos(-\pi-\alpha)]=0$, so $\tau_\alpha^vJ=1-v^2$ and
\[
k_1=\frac{\mathcal T}{w_1^3}-\TS_1,\qquad\mathcal T:=\Bigl(\frac{\tau_1}{\tau_\alpha}\Bigr)^v\ge1 .
\]
By Lemma~\ref{lem:transport}(a), $(1+v)w_1\ge\operatorname{Im}\Xi=\tau_1^v\sfc=(\cos\alpha-v)\mathcal T>0$ (\hyptag{SC2}), and $w_1\ge1$, so
$0\le\mathcal T/w_1^3\le(1+v)/(\cos\alpha-v)\le(1+v)/(\cos\frac12-v)$. In $\TS_1$, $|\sa\,\partial_\alpha\log C|\le1$, so by
Lemma~\ref{lem:transport}(e) $|\TS_1|\le\frac v{(1-v^2)\rho_1}[\rho_1+1+v]\le\frac{2v}{1-v^2}$. Hence $|k_1|\le k_{\mathrm{arm}}$.

\emph{Tilt.} By Lemma~\ref{lem:transport}(a), with $\sa=-\sin\alpha<0$ and $\sfc\tau_1^v=(\cos\alpha-v)\mathcal T>0$,
\[
(1-v^2)\,\partial_\alpha S_\alpha(t_1)=\sa[(1-v\cos\beta_1)-iv\sin\beta_1]+\sfc\tau_1^v[-v\sin\beta_1+i(1+v\cos\beta_1)]
\]
has negative real part and positive imaginary part; so the tangent is $e^{i(\pi/2+\delta)}$ with $\delta\in[0,\pi/2)$ and
\[
\tan\delta=\frac{\sin\alpha(1-v\cos\beta_1)+v\sin\beta_1(\cos\alpha-v)\mathcal T}{v\sin\alpha\sin\beta_1+(1+v\cos\beta_1)(\cos\alpha-v)\mathcal T}
\le\frac{(1+v)\sin\alpha}{(1-v)(\cos\alpha-v)\mathcal T}+\frac{v\sin\beta_1}{1-v};
\]
with $\mathcal T\ge1$, $\sin\alpha\le y$, $\sin\beta_1=y/\rho_1\le y/(1+v)$, $\cos\alpha\ge\cos\frac12$: $\tan\delta\le d_{\mathrm{arm}}y$.

\emph{Stretch.} $|\Xi|^2=\sin^2\!\alpha+(\cos\alpha-v)^2\mathcal T^2\le\mathcal T^2$, so $w_1\le\mathcal T/(1-v)$. By
Lemma~\ref{lem:coord}(d) and $\tau_1>1$: $(1-v)\tau_1^{1+v}\le2|h_1|/C+1+v$, with $|h_1|\le2-\cos\frac12$,
$C=\sg(\alpha)\ge2\tau_\alpha^{1-v}/(1+\tan^2\frac14)$ and $1+v\le(1+v)\tan(\frac14)^{1-v}\tau_\alpha^{-(1-v)}$ ($\tau_\alpha\le\tan\frac14$).
So $\tau_1^{1+v}\le\mathcal E_1\tau_\alpha^{-(1-v)}$, $\tau_1/\tau_\alpha\le\mathcal E_1^{1/(1+v)}\tau_\alpha^{-2/(1+v)}$ and
$\mathcal T\le\mathcal E_1^{v}(2/\alpha)^{2v}$ ($\mathcal E_1\ge1$, $\tau_\alpha\ge\alpha/2$, $2/\alpha>1$). Hence
$w_1\le W_{\mathrm{arm}}(2/\alpha)^{2v}$. Finally $y=C\tau_1^v\le2\tau_\alpha^{1-v}(\mathcal E_1\tau_\alpha^{-(1-v)})^{v/(1+v)}\le
2\mathcal E_1\tau_\alpha^{(1-v)/(1+v)}$.

(d) If $\circab\alpha\ge\frac12$, (b) with $\rho=\rho_{\mathrm{mid}}$ gives $|k_1|\le1+\vk(\rho_{\mathrm{mid}})$ and
$w_1\le W_1(\rho_{\mathrm{mid}})(\pi/2)^{2v}(2/\pi)^{2v}\le W_{\mathrm{all}}(2/\circab\alpha)^{2v}$. If $\circab\alpha<\frac12$, use (c) and
the mirror symmetry. $\rho_1\ge1+v$: the argument of (c) uses only $h_1\le-(1-v^2)$, which holds for all $\alpha$ (for $\alpha=\pi$,
$\beta\equiv\pi$ and $\rho_1=1+(1+v)t_1$).

(e) By Lemma~\ref{lem:coord}(c), on a straight run $\rho\ge C/\gmax>0$ whenever $0<\beta<\pi$ and $C>0$. In run 1,
$C=\sg(\alpha)>0$ for $\alpha\in(0,\pi)$ (mirror for $(\pi,2\pi)$), and the label $\pi$ has $\rho\ge1$. In run 2,
$\sbb\in(0,\pi)$ for $\alpha\ne\alpha_0$ (Lemma~\ref{lem:tang}(a); for $\alpha=\pi$, $\sbb=\pi-\beta^*$), so $C=\sr\sg(\sbb)>0$; and the
label $\alpha_0$ has $\sbb=0$, $\rho=\rho_0-(1-v)(t-t_1)>0$ for $t<\tc$, while $\ttd<\tc$ (Lemma~\ref{lem:tang}(b), whose proof does not
use (e)). The label $0$ lies on the axis ahead of $P$: $\rho(0)=1$, $\rho'=-(1-v)$; it is caught, a fortiori seen, at $1/(1-v)$.
\end{proof}

\begin{lemma}[tangency; data at $t_1$]\label{lem:tang}
\begin{enumerate}
\item[(a)] For $\theta\in(0,\pi)$: $\rho_1\,\partial_\theta\beta_1=\sin\beta_1\cdot t_1(\cos\alpha_0-\cos\theta)/\sin\theta$. Hence $\beta_1$ has its
unique minimum on $(0,\pi)$ at $\alpha_0$; there $\sbb=\beta_1-\beta^*\in[0,\pi/2)$; on the lower half $\sbb\in(0,\pi)$; and
\[
\partial_\theta\sbb=\Lambda(\theta)\sin\frac{\theta-\alpha_0}2,\qquad
\Lambda(\theta):=\frac{2t_1\sin\beta_1\sin((\theta+\alpha_0)/2)}{\rho_1\sin\theta}>0\qquad(0<\theta<\pi).
\]
So $\sbb$ has a double zero at $\alpha_0$: the true front at $t_1$ is tangent at $S_{\alpha_0}(t_1)$ to the ray of run 2.
\item[(b)] $\alpha_0\in[\underline\alpha,\pi/3]$; $|\rho_0-\sqrt3|\le e_\rho$; $\rho_0>\eps$; $\ttd\le\tc\le T_2$,
$0<\tc-\ttd\le v\eps/(3(1-v)^2)$ and $\ttd-\ta\le\eps/(1-v)$; every candidate, true or virtual, satisfies
$|x_\theta(t)-e^{i\theta}|\le vt\le d_{\mathrm s}$ for $t\le\ttd$; $P([t_1,\ttd])$ lies within $d_P$ of the segment $[2,e^{i\pi/3}]$; the
virtual candidate $\alpha_0$ stays on the ray of run 2, $\rho_{\alpha_0}(\ta)=\eps$ and $P(\ttd)=x_{\alpha_0}(\ttd)$.
\item[(c)] $I_\eta\subset[\theta_L,\theta_R]\subset[\frac12,\frac\pi2]$, and on $I_\eta$:
\begin{gather*}
\sr_{\min}\le\sr\le\sr_{\max},\quad\Lambda_{\min}\le\Lambda\le\Lambda_{\max},\quad\ell_{\min}\le\sr\Lambda\le\ell_{\max},\quad
\sA_{\min}\le\sA\le\sA_{\max},\\
1\le w_1\le W_{\mathrm w},\quad|k_1-1|\le\vk_{\mathrm w},\quad b_{\mathrm{lo}}\xi^2\le\sbb\le b_{\mathrm{hi}}\xi^2\le\bar b,\quad
\sbb/2\le\tau_{\sbb}\le\beta_2\xi^2,\\
C_{\mathrm{lo}}\xi^{2-2v}\le C\le C_{\mathrm{hi}}\xi^{2-2v},\quad\sa_{\min}\le\sa\le\sa_{\max},\quad|\partial_\theta\log C|\le L_C/\xi;
\end{gather*}
for $\theta<\alpha_0$: $\partial_\theta\log C\le-(L_{\mathrm{lo}}-\eta\sA_{\max}/\sr_{\min})/\xi<0$, so $C$ is strictly monotone in $\xi$ on
that side; and the quantity $\TS$ of Lemma~\ref{lem:transport} at $t_1$, in the frame of run 2, satisfies $|\TS|\le S_1v$.
\item[(d)] For $\theta\in(0,\pi)\setminus\{\alpha_0\}$, at $t_1$ in the frame of run 2,
\[
\tau_{\sbb}^v\,\partial_\theta C=\frac{2\tau_{\sbb}\sA}{1+\tau_{\sbb}^2}+\sr\,\partial_\theta\sbb\,(\cos\sbb-v).
\]
Consequently, on $I_\eta$: for $\theta>\alpha_0$, $0<\tau_{\sbb}^v\partial_\theta C\le\xi(\ell_{\max}/2+\sA_{\max}b_{\mathrm{hi}}(1+e_\tau)\eta)$;
for $\theta<\alpha_0$, $\partial_\theta C<0$ and $|\tau_{\sbb}^v\partial_\theta C|\le\sr|\partial_\theta\sbb|\le\ell_{\max}\xi/2$.
\end{enumerate}
\end{lemma}
\begin{proof}
(a) From $C=\rho\sg(\beta)$, $h=\rho(\cos\beta+v)$ and $\sg'/\sg=(\cos\beta-v)/\sin\beta$, on any straight run
\[
(\cos\beta+v)\,\partial_\theta C-\sg\,\partial_\theta h=\rho\,\partial_\theta\beta\cdot\sg(1-v^2)/\sin\beta .
\]
In run 1 at $t_1$: $\partial_\theta C=\sg(\theta)(\cos\theta-v)/\sin\theta$, $\partial_\theta h=-\sin\theta$, $\sg(\theta)/\sg(\beta_1)=\rho_1$ and
$\rho_1(\cos\beta_1+v)=h_1$; so $(1-v^2)\rho_1\partial_\theta\beta_1=\sin\beta_1[h_1(\cos\theta-v)+\sin^2\!\theta]/\sin\theta$, and
$h_1(\cos\theta-v)+\sin^2\!\theta=(1-v^2)(1-t_1(\cos\theta-v))=(1-v^2)t_1(\cos\alpha_0-\cos\theta)$, as $t_1(\cos\alpha_0-v)=1$. Since
$h_1<0$, $\beta_1\in(\arccos(-v),\pi)$, so $\beta^*>\pi/2$ and on the upper half $\sbb=\beta_1-\beta^*\in[0,\pi/2)$. The mirror point
of $S_\alpha$ has polar angle $-\beta_1(\alpha)$ about the axis of run 1, i.e.\ $\sbb=2\pi-\beta_1(\alpha)-\beta^*\in(\pi-\beta^*,2\pi-2\beta^*]\subset(0,\pi)$.
The formula for $\partial_\theta\sbb$ is $\cos\alpha_0-\cos\theta=2\sin\frac{\theta+\alpha_0}2\sin\frac{\theta-\alpha_0}2$. In the frame of run 2,
$\partial_\theta(S_\theta-P)(t_1)=(\sA+i\sr\,\partial_\theta\sbb)e^{i\sbb}$, which is real at $\alpha_0$.

(b) $t_1=(2-v)/(1-v)$, so $\cos\alpha_0-\frac12=v(3-2v)/(2(2-v))\in[0,v]$ and $\alpha_0\in[\underline\alpha,\pi/3]$. By Lemma~\ref{lem:run1}(a),
$|S_{\alpha_0}(t_1)-P(t_1)-(e^{i\pi/3}-2)|\le vt_1+|e^{i\alpha_0}-e^{i\pi/3}|+|t_1-2|\le vt_1+\Delta_{\mathrm a}+\mu\le e_\rho$, and
$|e^{i\pi/3}-2|=\sqrt3$; so $|\rho_0-\sqrt3|\le e_\rho$, $\rho_0\ge\sqrt3-e_\rho>\eps$ (\hyptag{SC1}), and $\tc=t_1+\rho_0/(1-v)\le T_2$.
The true candidate $\alpha_0$ has $\sbb=0$: it stays on the ray, at distance $\rho_0-(1-v)(t-t_1)$ (Proposition~\ref{prop:runs}), which
is $\eps$ at $\ta$; up to $\ta$ it is the virtual one (Lemma~\ref{lem:truefront}). For $t\ge\ta$, $x=P+\sigma\nhat$ solves
$x'=vF(x-P)$ iff $\sigma'=-1+vf(\sigma)$, $\sigma(\ta)=\eps$; by uniqueness this is the virtual candidate $\alpha_0$, and $\sigma$ reaches
$0$ after the time $\int_0^\eps d\sigma/(1-vf(\sigma))=\eps\int_0^1d\sigma/(1-v\sigma(2-\sigma))=\ttd-\ta\le\eps/(1-v)$. Hence
$P(\ttd)=x_{\alpha_0}(\ttd)$, and, as $f<1$ on $[0,\eps)$ and $\int_0^\eps(1-f)=\eps/3$,
$\tc-\ttd=\int_0^\eps\frac{v(1-f)}{(1-v)(1-vf)}\,d\sigma\in(0,v\eps/(3(1-v)^2)]$. Positions: speed $\le v$ and $t\le\ttd\le T_2$. Finally
$P(t_1)=2+\mu$ is within $\mu$ of the segment, $P(\ttd)=x_{\alpha_0}(\ttd)$ is within $vT_2+\Delta_{\mathrm a}\le d_{\mathrm s}+\Delta_{\mathrm a}$
of $e^{i\pi/3}$, and the distance to a convex set is convex along the segment $P([t_1,\ttd])$.

(c) $I_\eta\subset[\theta_L,\theta_R]$ by (b), $\subset[\frac12,\frac\pi2]$ by \hyptag{SC3}, \hyptag{SC4}. Window labels are $\ge\theta_L>\as$,
and $\rho_\theta\ge\sin\theta\ge\rho_{\mathrm w}$ on $[0,t_1]$, so Lemma~\ref{lem:run1}(b) with $\rho=\rho_{\mathrm w}$ gives the bounds on
$w_1$, $k_1$, and $\partial_\theta S_\theta(t_1)=w_1e^{i\psi_1}$ with $|\psi_1-(\theta+\pi/2)|\le\omega_{\mathrm w}$.

\emph{$\sr$.} $\mathsf p:=S_\theta(t_1)-P(t_1)=e^{i\theta}-2+\mathsf d$ with $|\mathsf d|\le vt_1+\mu\le d_1$, and $R_0$ is increasing on $(0,\pi)$.

\emph{$\Lambda$.} $\sin\beta_1/\rho_1=\operatorname{Im}\mathsf p/\sr^2$ with $\operatorname{Im}\mathsf p=\operatorname{Im}S_\theta(t_1)\in[\sin\theta,\sin\theta+vt_1]$
($y$ non-decreasing in run 1, speed $\le v$); so $\Lambda=2t_1(\operatorname{Im}\mathsf p/\sin\theta)\sin((\theta+\alpha_0)/2)/\sr^2$ with
$1\le\operatorname{Im}\mathsf p/\sin\theta\le\iota_{\mathrm w}$, $t_1\in[2,\bar t_1]$ and
$(\theta+\alpha_0)/2\in[\underline\alpha-\eta/2,\pi/3+\eta/2]\subset(0,\pi/2]$ (\hyptag{SC3}, \hyptag{SC5}), where $\sin$ is increasing.

\emph{$\sr\Lambda$} $=2t_1(\operatorname{Im}\mathsf p/\sin\theta)\,\tilde q(\theta)\,(R_0/\sr)$ with $\tilde q:=\sin((\theta+\alpha_0)/2)/R_0$. Here
$\frac d{d\theta}\log\tilde q=\frac12\cot\frac{\theta+\alpha_0}2-s(\theta)$, $s=R_0'/R_0$; $s'$ has the sign of $5\cos\theta-4<0$
(\hyptag{SC6}), so $s\ge s(\theta_R)$, while $\frac12\cot\frac{\theta+\alpha_0}2\le\frac12\cot(\underline\alpha-\eta/2)$; by \hyptag{SC7},
$\tilde q$ is decreasing on $I_\eta$. Hence $\tilde q\in[\tilde q(\alpha_0+\eta),\tilde q(\alpha_0-\eta)]\subset[q_{\min},q_{\max}]$, and
$R_0/\sr\in[1/(1+d_1/R_L),1/(1-d_1/R_L)]$ (\hyptag{SC8}).

\emph{$\sA$} $=\partial_\theta S_\theta(t_1)\cdot\hat{\mathsf p}=w_1\cos\angle(e^{i\psi_1},\hat{\mathsf p})$. With $\mathsf p_0:=e^{i\theta}-2$:
$ie^{i\theta}\cdot\mathsf p_0=2\sin\theta>0$ and $ie^{i\theta}\times\mathsf p_0=2\cos\theta-1$, so
$\angle(ie^{i\theta},\mathsf p_0)=\arcsin(|2\cos\theta-1|/R_0)\le\arcsin m_0$; and $\angle(\mathsf p_0,\mathsf p)\le\arcsin(|\mathsf d|/|\mathsf p_0|)\le\omega_{\mathrm p}$.
So $\angle(e^{i\psi_1},\hat{\mathsf p})\le\angle_{\mathrm w}<\pi/2$ (\hyptag{SC9}) and $\sA\in[\cos\angle_{\mathrm w},W_{\mathrm w}]$.

\emph{$\sbb$} $(\alpha_0\pm\xi)=\int_0^\xi\Lambda(\alpha_0\pm\xi')\sin(\xi'/2)\,d\xi'\in[\Lambda_{\min},\Lambda_{\max}]\cdot4\sin^2(\xi/4)$, and
$(\xi^2/4)(1-\xi^2/48)\le4\sin^2(\xi/4)\le\xi^2/4$. Since $\tan s/s$ is increasing on $(0,\pi/2)$ and $\sbb\le\bar b<\pi/2$ (\hyptag{SC10}):
$\sbb/2\le\tau_{\sbb}\le(\sbb/2)(1+e_\tau)\le\beta_2\xi^2$.

\emph{$C$} $=2\sr\tau_{\sbb}^{1-v}/(1+\tau_{\sbb}^2)$: $C\le2\sr_{\max}(\beta_2\xi^2)^{1-v}\le C_{\mathrm{hi}}\xi^{2-2v}$ (\hyptag{SC11}), and
$C\ge2\sr_{\min}(b_{\mathrm{lo}}\xi^2/2)^{1-v}/(1+\tan^2(\bar b/2))\ge C_{\mathrm{lo}}\xi^{2-2v}$ (\hyptag{SC12}).

\emph{$\sa$} $=\sA(\cos\sbb+v)-\sr\sin\sbb\,\partial_\theta\sbb$, and $|\sr\sin\sbb\,\partial_\theta\sbb|\le\bar b\,\ell_{\max}\xi/2\le\mathrm{aux}$, as
$|\partial_\theta\sbb|\le\Lambda\sin(\xi/2)\le\Lambda\xi/2$.

\emph{$\partial_\theta\log C$} $=\sA/\sr+(\cos\sbb-v)\,\partial_\theta\sbb/\sin\sbb$. Now
\[
\frac{|\partial_\theta\sbb|}{\sbb}=\frac{\Lambda(\theta)\sin(\xi/2)}{\int_0^\xi\Lambda\sin(\xi'/2)d\xi'}\in
\Bigl[\frac{\Lambda_{\min}}{\Lambda_{\max}},\frac{\Lambda_{\max}}{\Lambda_{\min}}\Bigr]\cdot\frac12\cot\frac\xi4,
\]
and for $0<y\le1$, $1/y-y/2\le\cot y\le1/y$ (the function $y\cos y-\sin y+\frac{y^2}2\sin y$ vanishes at $0$ and has derivative
$\frac{y^2}2\cos y\ge0$), so $(2/\xi)(1-\xi^2/32)\le\frac12\cot(\xi/4)\le2/\xi$. Moreover
$(\cos\sbb-v)\sbb/\sin\sbb\in[\cos\bar b-v,\sbb\cot\sbb]\subset[\cos\bar b-v,1]$, and $\cos\bar b-v>0$ (\hyptag{SC14}). Hence
$(\cos\sbb-v)|\partial_\theta\sbb|/\sin\sbb\in[L_{\mathrm{lo}}/\xi,\,2\Lambda_{\max}/(\Lambda_{\min}\xi)]$, while
$0<\sA/\sr\le\sA_{\max}/\sr_{\min}\le(\eta\sA_{\max}/\sr_{\min})/\xi$. For $\theta>\alpha_0$ both terms are positive and
$\partial_\theta\log C\le L_C/\xi$; for $\theta<\alpha_0$, $\partial_\theta\log C\le\sA_{\max}/\sr_{\min}-L_{\mathrm{lo}}/\xi\le
-(L_{\mathrm{lo}}-\eta\sA_{\max}/\sr_{\min})/\xi<0$ (\hyptag{SC15}), and $\ge-2\Lambda_{\max}/(\Lambda_{\min}\xi)\ge-L_C/\xi$.

\emph{$\TS$ at $t_1$.} The tangent is the direction of $(\sA+i\sr\,\partial_\theta\sbb)e^{i\sbb}$, so $|\sin\phi|=\sr|\partial_\theta\sbb|/w_1\le\ell_{\max}\xi/2$;
by Lemma~\ref{lem:transport}(e) with $\rho=\sr\ge\sr_{\min}$ and $w_1\ge1$:
$|\TS|\le\frac{v\ell_{\max}\xi}{2(1-v^2)\sr_{\min}}[\sr_{\max}\sa_{\max}L_C/\xi+1+v]\le S_1v$ ($\xi\le\eta$).

(d) $\partial_\theta C=\sg(\sbb)\sA+\sr\sg'(\sbb)\partial_\theta\sbb$ with $\sg'=\sg(\cos\beta-v)/\sin\beta$, and $\tau^v\sg(\beta)=\sin\beta=2\tau/(1+\tau^2)$
(Lemma~\ref{lem:coord}(a)); this is the identity. On the forward side both terms are positive,
$2\tau_{\sbb}\sA/(1+\tau_{\sbb}^2)\le2\beta_2\xi^2\sA_{\max}=\sA_{\max}b_{\mathrm{hi}}(1+e_\tau)\xi^2\le\sA_{\max}b_{\mathrm{hi}}(1+e_\tau)\eta\,\xi$, and
$0<\sr\,\partial_\theta\sbb(\cos\sbb-v)\le\sr\Lambda\xi/2\le\ell_{\max}\xi/2$. On the tail side $\partial_\theta\log C<0$ by (c), so
$|\tau_{\sbb}^v\partial_\theta C|=\sr|\partial_\theta\sbb|(\cos\sbb-v)-2\tau_{\sbb}\sA/(1+\tau_{\sbb}^2)\le\sr|\partial_\theta\sbb|$.
\end{proof}

\begin{lemma}[forward arc]\label{lem:fwd}
Let $\theta=\alpha_0+\xi$, $0<\xi\le\eta$. For every $t\in[t_1,\tc)$, in the frame of run 2:
\begin{enumerate}
\item[(a)] $h(t)\ge h_{\mathrm f}\xi+(1-v^2)\omega(t)\ge h_{\mathrm f}(\xi+\omega(t))>0$ and $\rho\ge h/(1+v)\ge\rho_{\mathrm c}(\xi+\omega)$;
in particular $\rho>0$ and $0<\beta<\pi$.
\item[(b)] $\tau\le\tau_{\sbb}((1+v)\sr/h)^{1+\mu}$, $(\tau/\tau_{\sbb})^v\le(Q/(\xi+\omega))^\mu$, and
$0<y\le y_{\mathrm f}\xi^2(Q/(\xi+\omega))^\mu$.
\item[(c)] $0<Y_r\le Y_1\xi(Q/(\xi+\omega))^\mu\le\bar Y$; $|\psi|\le|\tan\psi|\le\theta_{\mathrm f}\xi(Q/(\xi+\omega))^\mu\le
\theta_{\mathrm f}Q^{\mub}\xi^{1-\mu}$; $1\le w\le W_{\mathrm f}$.
\item[(d)] At $\ttd$: $\tau\le\tau_{\mathrm f}\xi^{1-\mu}\le\bar\tau_{\mathrm f}$, and $\rho(t)\ge\rho(\ttd)\ge\rho_{\mathrm c}\xi$ for all
$t\in[t_1,\ttd]$; $\rho\ge\eps$ on $[0,t_1]$. Hence every label with $\rho_{\mathrm c}\xi>\eps$, in particular every
$\xi\ge\xi_1\eps$ and a neighbourhood of $\xi=\xi_1\eps$ (\hyptag{SC17}), has $\rho>\eps$ on $[t_1,\ttd]$: it is true up to $\ttd$,
virtual $=$ true there, and the front is $C^\infty$ near it. For $\xi_1\eps\le\xi\le\eta$ moreover $\rho(\ttd)\le\rho_{\mathrm u}\xi$
and $|k|\le k_{\mathrm f}\xi^{-\mu}$ at $\ttd$.
\item[(e)] The true front at every $t\in[t_1,\tc)$ is tangent to $\nhat$ at the candidate $\alpha_0$.
\end{enumerate}
\end{lemma}
\begin{proof}
(a) $\sr-\rho_0=\int_{\alpha_0}^\theta\sA\ge\sA_{\min}\xi$ and $\sr(1-\cos\sbb)\le\sr_{\max}\sbb^2/2\le(\sr_{\max}b_{\mathrm{hi}}^2\eta^3/2)\xi$
(Lemma~\ref{lem:tang}(c)); insert into \eqref{eq:h2} and use $h_{\mathrm f}\le1-\vbar^2\le1-v^2$ (\hyptag{SC16}). $h=\rho(\cos\beta+v)\le(1+v)\rho$.
$\rho>0$ and Proposition~\ref{prop:runs} give $0<\beta<\pi$ (as $0<\sbb<\pi$ at $t_1$).

(b) \eqref{eq:tau2}, $v(1+\mu)=\mu$, $(1+v)\sr/h\le Q/(\xi+\omega)$ by (a), $\tau_{\sbb}\le\beta_2\xi^2$ and $\sr\le\sr_{\max}$.

(c) By Lemma~\ref{lem:transport}(b) and Lemma~\ref{lem:tang}(d), $Y_r=(\tau/\tau_{\sbb})^v\tau_{\sbb}^v\partial_\theta C/\sa\in(0,Y_1\xi(Q/(\xi+\omega))^\mu]$
($\sa\ge\sa_{\min}>0$, \hyptag{SC13}), and $Y_1\xi(Q/(\xi+\omega))^\mu\le Y_1Q^\mu\xi^{1-\mu}\le\bar Y$ (\hyptag{SC18}). So
$N_\psi\in[-v\sin\beta,(1+v)Y_r]$ and $D_\psi\ge1-v-v\bar Y>0$ (\hyptag{SC19}); hence $|\psi|\le|\tan\psi|$. With
$\sin\beta=y/\rho\le(y_{\mathrm f}/\rho_{\mathrm c})\xi(Q/(\xi+\omega))^\mu$:
$|\tan\psi|\le\max((1+v)Y_1,vy_{\mathrm f}/\rho_{\mathrm c})\,\xi(Q/(\xi+\omega))^\mu/(1-v-v\bar Y)\le\theta_{\mathrm f}\xi(Q/(\xi+\omega))^\mu$.
$w=\sa|D_\psi+iN_\psi|/(1-v^2)\le W_{\mathrm f}$, since $|D_\psi|\le1+v+v\bar Y$, $|N_\psi|\le(1+v)\bar Y+v$; $w\ge1$ by Lemma~\ref{lem:truefront}.

(d) By (b) with $h\ge h_{\mathrm f}\xi$: $\tau(\ttd)\le\beta_2\xi^2(Q/\xi)^{1+\mu}\le\tau_{\mathrm f}\xi^{1-\mu}\le\bar\tau_{\mathrm f}$, so
$\cos\beta\ge1-2\bar\tau_{\mathrm f}^2>v$ (\hyptag{SC20}). As $\tau$ increases in time, $\cos\beta>v$ on $[t_1,\ttd]$ and $\rho$ is
decreasing there (Lemma~\ref{lem:coord}(c)): $\rho(t)\ge\rho(\ttd)\ge\rho_{\mathrm c}\xi$. In run 1, $\rho\ge\eps$ since
$\theta\ge\theta_L>\as$ (Lemma~\ref{lem:run1}(a)). Virtual $=$ true by Lemma~\ref{lem:truefront}. For $\xi\ge\xi_1\eps$:
$\omega(\ttd)\le v\eps/(3(1-v)^2)\le\delta_{\mathrm{fw}}\xi$, so by \eqref{eq:h2} $h\le(1+v)\sA_{\max}\xi+\delta_{\mathrm{fw}}\xi$, and
$\rho=h/(\cos\beta+v)\le h/(1-2\bar\tau_{\mathrm f}^2)\le\rho_{\mathrm u}\xi$.

\emph{Curvature.} Lemma~\ref{lem:transport}(d) with $t'=t_1$, $t''=\ttd$: $(\tau/\tau_{\sbb})^v\le(Q/\xi)^\mu\le Q^{\mub}\xi^{-\mu}$,
$|k_1|\le1+\vk_{\mathrm w}$ and $|\TS'|\le S_1v$ (Lemma~\ref{lem:tang}(c)). For $\TS''$, Lemma~\ref{lem:transport}(e) with
$|\sin\phi|=|\sin(\beta-\psi)|\le\sin\beta+|\psi|\le2\tau+|\psi|\le(2\tau_{\mathrm f}+\theta_{\mathrm f}Q^{\mub})\xi^{1-\mu}$,
$\rho\ge\rho_{\mathrm c}\xi$ and $\rho|\sa||\partial_\theta\log C|/w^2\le\rho_{\mathrm u}\xi\,\sa_{\max}L_C/\xi$ gives $|\TS''|\le S_2v\xi^{-\mu}$.
So $|k|\le Q^{\mub}\xi^{-\mu}(1+\vk_{\mathrm w}+S_1v)+S_2v\xi^{-\mu}\le k_{\mathrm f}\xi^{-\mu}$.

(e) For $t<\tc$ the true front is smooth near the label $\alpha_0$ ($\rho_{\alpha_0}(t)>0$), and by (c) $\psi\to0$ as $\xi\downarrow0$ with
$\omega(t)>0$ fixed.
\end{proof}

\begin{lemma}[virtual labels next to $X_0$]\label{lem:virt}
At $\ttd$, for every $\theta=\alpha_0+\xi$, $0\le\xi\le\eta$, the virtual front has tangent $e^{i\psi}\nhat$ with
$|\psi(\theta,\ttd)|\le\theta_{\mathrm v}\xi^{1-\mu}$, and $1\le w\le W_{\mathrm v}$; $\psi(\alpha_0,\ttd)=0$. For $\xi\ge\xi_1\eps$ the
virtual front is the true one.
\end{lemma}
\begin{proof}
For $\xi_1\eps\le\xi\le\eta$: Lemma~\ref{lem:fwd}(c),(d) at $\ttd$, as $\theta_{\mathrm f}Q^{\mub}\le\theta_{\mathrm v}$ and
$W_{\mathrm f}\le W_{\mathrm v}$. Let $0<\xi<\xi_1\eps$; then $\xi<\eta$, since $\xi_1\eps\le u_{\mathrm v}\epsbar\le\eta/2$ (\hyptag{SC40},
$W_{\mathrm v}\ge1$). On $[0,t_1]$ the label is true with $\rho\ge\eps$ (Lemma~\ref{lem:run1}(a)). On $[t_1,\ta]$, by \eqref{eq:h2},
$h_\theta(t)-h_{\alpha_0}(t)=(1+v)(\sr-\rho_0)-\sr(1-\cos\sbb)\ge h_{\mathrm f}\xi>0$ (proof of Lemma~\ref{lem:fwd}(a)) and
$h_{\alpha_0}(t)=(1+v)\rho_{\alpha_0}(t)\ge(1+v)\eps$ (the candidate $\alpha_0$ is on the ray), so $(1+v)\rho_\theta\ge h_\theta>(1+v)\eps$.
Hence virtual $=$ true on $[0,\ta]$ for all labels in $(\alpha_0,\alpha_0+\xi_1\eps)$ (Lemma~\ref{lem:truefront}), and by
Lemma~\ref{lem:fwd}(b),(c) at $\ta$ ($\omega(\ta)=\eps/(1-v)\ge\eps$):
$|\psi(\ta)|\le\theta_{\mathrm f}\xi(Q/\xi)^\mu$, $0\le y(\ta)\le y_{\mathrm f}\xi^2(Q/\xi)^\mu$, $w(\ta)\le W_{\mathrm f}$. On
$[\ta,\ttd]$, of length $I\le\eps/(1-v)$ (so $vI/\eps\le\mu$), the virtual field acts. With $r=x_\theta-P=(z+iy)\nhat$:
$r'=vF(r)-\nhat$, so $y'=v(f(\rho)/\rho)y$ with $0\le f(\rho)/\rho\le2/\eps$, hence $0\le y(t)\le y(\ta)e^{2\mu}$. By
Lemma~\ref{lem:flow}(b), $|\psi_t|\le v\rho|\sin\phi|/\max(\rho,\eps)^2$, where $\rho|\sin\phi|=|T\times r|=|y\cos\psi-z\sin\psi|\le y+\rho|\psi|$
and $\rho/\max(\rho,\eps)^2\le1/\eps$; so $|\psi|'\le(v/\eps)|\psi|+vy/\eps^2$ almost everywhere. By Gronwall,
\[
|\psi(\ttd)|\le e^{\mu}\Bigl[|\psi(\ta)|+\frac{vI}{\eps^2}e^{2\mu}y(\ta)\Bigr]\le e^\mu\bigl[\theta_{\mathrm f}+\mu e^{2\mu}y_{\mathrm f}\,\xi/\eps\bigr]
Q^\mu\xi^{1-\mu}\le\theta_{\mathrm v}\xi^{1-\mu}
\]
($\xi/\eps<\xi_1$). By Lemma~\ref{lem:flow}(b) and Lemma~\ref{lem:field}(ii), $w_t\le va^*w\le(2v/\eps)w$, so
$w(\ttd)\le w(\ta)e^{2\mu}\le W_{\mathrm v}$; $w\ge1$ by Lemma~\ref{lem:flow}(a). For $\xi=0$: the virtual front is $C^1$ and coincides
at $\ta$ with the true front on $(\alpha_0,\alpha_0+\xi_1\eps)$, so $w(\alpha_0,\ta)\le W_{\mathrm f}$ and its tangent at $\alpha_0$ is the
limit of the true tangents, $\nhat$ (Lemma~\ref{lem:fwd}(e)); thus $\psi(\ta)=0=y(\ta)$, and the same inequalities give $\psi(\alpha_0,\ttd)=0$
and $w(\alpha_0,\ttd)\le W_{\mathrm v}$.
\end{proof}
